\documentclass[10pt]{article}
\usepackage[margin=1.35cm]{geometry}
\usepackage{graphicx}
\usepackage{amsmath}
\usepackage{booktabs}
\usepackage{xcolor}
\pagestyle{empty}
\setlength{\parindent}{0pt}
\setlength{\parskip}{3pt}
\graphicspath{{../analysis_output/accessibility_alpha_v2w4wd/}}

\begin{document}
\begin{center}
{\Large\bfseries Preliminary matched-$V_2$/readout-decay experiment}\par\vspace{-1pt}
{\small Seed 0; 8 ReLU scales $\times$ 8 decay values = 64 completed cells}
\end{center}
\vspace{-4pt}

\textbf{Question and design.}
This experiment repeats the unfrozen accessibility sweep while applying the same phase-2
weight decay $\lambda$ to $V_2$ and to the $d4$ readout row $W[2]$.  The purpose is to remove
the V2-only experiment's scale escape, in which $V_2$ shrank while the unpenalized readout
grew.  The preliminary grid covers
$\alpha\in\{0.1,0.3,0.5,0.75,1,1.5,2.5,4\}$ and
$\lambda\in\{0,10^{-4},3\!\cdot\!10^{-4},10^{-3},3\!\cdot\!10^{-3},10^{-2},
3\!\cdot\!10^{-2},10^{-1}\}$.

\textbf{Is $R_0$ controlled by the ReLU coefficient?  Yes.}
At $d4$ onset the pretrained first-block matrix is replaced by $V_1\mapsto\alpha V_1$.
For $\alpha>0$, ReLU homogeneity gives
$\operatorname{ReLU}(h_0(\alpha V_1)^\top)=
\alpha\operatorname{ReLU}(h_0V_1^\top)$, so $\alpha$ directly controls the first-block
features and therefore the measured accessibility $R_0=Q_2/Q_1$.  The relation is monotone
but not linear because the residual term $h_0$ is not rescaled.  In seed 0, each $\alpha$
gave exactly one $R_0$ across all eight $\lambda$ values:

\vspace{-2pt}
\begin{center}\small
\begin{tabular}{@{}c rrrrrrrr@{}}
\toprule
$\alpha$ & 0.10 & 0.30 & 0.50 & 0.75 & 1.00 & 1.50 & 2.50 & 4.00 \\
$R_0$    & .038 & .113 & .181 & .256 & .319 & .411 & .510 & .568 \\
\bottomrule
\end{tabular}
\end{center}
\vspace{-2pt}
Thus the theoretical boundary $R_c=1/\sqrt{8}=0.354$ is crossed between
$\alpha=1$ and $\alpha=1.5$.  Across seeds, $\alpha$ remains the controlled intervention and
$R_0$ the measured response; small seed-to-seed shifts are expected.

\textbf{Preliminary result.}
Matched readout decay does \emph{not} sharpen a compositional transition near $R_c$ in seed 0.
The largest compositional shares occur only under very weak decay: for the supercritical
$\alpha=2.5$ cell, $f_{2+2}$ rises from 0.085 at $\lambda=0$ to 0.114 at
$3\!\cdot\!10^{-4}$, then collapses to 0.017 at $10^{-3}$ and approximately zero by
$3\!\cdot\!10^{-3}$.  Averaged across $\alpha$, $f_{2+2}$ falls from 0.043 at $\lambda=0$
to 0.008 at $10^{-3}$, 0.002 at $10^{-2}$, and 0.0005 at $10^{-1}$.

\begin{figure}[h!]
\centering
\includegraphics[width=0.99\textwidth]{fig1_seed0_preliminary.pdf}
\vspace{-7pt}
\caption{Direct seed-0 analogue of Fig.~1 in the V2-only report.  Color is the weighted
$2{+}2\!\to\!4$ fraction; the dashed line is $R_c$.  Despite upward/rightward drift in
$R_{\rm final}$ under decay, cells remain blue: becoming accessible does not produce a
compositional phase when the $d4$ readout is decayed together with $V_2$.}
\end{figure}
\vspace{-7pt}

\textbf{Interpretation.}
The scale escape is removed---both $V_2$ and its $d4$ readout are penalized---but the network
responds by suppressing the measured block-2 compositional route rather than balancing it.
At the same time, $R(t)$ generally drifts upward (for example, at $\alpha=0.1$,
$R_{\rm final}=0.139$ at $\lambda=0$ and 0.368 at $10^{-2}$).  The coexistence of high
$R_{\rm final}$, near-zero $f_{2+2}$, and solved $d4$ is consistent with the first-block $d4$
leak/residual shortcut taking over.  This is a strong preliminary pattern, but uncertainty and
cross-seed consistency must wait for the full 192-cell merge.

\end{document}
